\begin{filecontents*}{\jobname.bib}
@article{huyen2024ai,
  author  = {Huyen, Chip},
  title   = {AI Engineering: Building Applications with Foundation Models},
  journal = {O'Reilly Media},
  year    = {2024}
}

@article{kottwitz2021latex,
  author  = {Kottwitz, Stefan},
  title   = {LaTeX Beginner's Guide: Create High-Quality Books, Articles, and Presentations},
  journal = {Packt Publishing},
  year    = {2021}
}
\end{filecontents*}

\documentclass[12pt,a4paper]{article}

% ---------- Packages ----------
\usepackage[utf8]{inputenc}
\usepackage[T1]{fontenc}
\usepackage[margin=2.54cm]{geometry}
\usepackage{booktabs,tabularx}
\usepackage{pgfgantt}
\usepackage[dvipsnames,table]{xcolor}
\usepackage[hidelinks]{hyperref}

% =========================================================================
% [TODO - TASK 1]: MACRO & METADATA CONFIGURATION
% Architectural Why: Hardcoding metadata across multiple chapters, headers,
% or declarations violates DRY principles and causes drift. Preamble macros
% act as a single source of truth for your dissertation metadata.
% Action: Replace the placeholders below with your actual project details.
% =========================================================================
\newcommand{\studentname}{[TODO: Enter Your Full Name]}
\newcommand{\studentid}{[TODO: Enter Coventry/TKH Student ID]}
\newcommand{\projecttitle}{[TODO: Enter Working FYP Project Title]}
\newcommand{\modulecode}{KH6002CMD}
\newcommand{\moduletitle}{Project Discovery}

% Document Header Styling
\title{\vspace{-1.5cm}\textbf{\large \modulecode: \moduletitle}\\[0.3cm]
\Large \textbf{\projecttitle}\\[0.2cm]
\normalsize Guided Practice Sandbox: LaTeX, BibTeX \& Project Planning}
\author{\textbf{\studentname} \quad (ID: \studentid)}
\date{\today}

\begin{document}
\maketitle
\thispagestyle{empty}

\noindent\rule{\textwidth}{0.8pt}
\vspace{-0.3cm}

% =========================================================================
% SECTION 1: CITATION PIPELINE & THE [?] TRAP
% =========================================================================
\section{Academic Engine \& Citation Pipeline}

Academic writing requires verifiable claims backed by peer-reviewed literature. LaTeX isolates bibliographic metadata in \texttt{.bib} databases and relies on a multi-pass compilation sequence (\texttt{pdflatex} $\rightarrow$ \texttt{bibtex} $\rightarrow$ \texttt{pdflatex} $\rightarrow$ \texttt{pdflatex}) to build numerical citations.

% =========================================================================
% [TODO - TASK 2]: THE [?] UNDEFINED CITATION TRAP
% Architectural Why: pdflatex does not parse .bib files directly; it logs
% citation requests into \jobname.aux. If a citation key is misspelled or
% not yet indexed by bibtex, pdflatex inserts a bold '[?]' into the PDF.
% Action:
% 1. Notice the two citations below: 'huyen2099_wrong_key' and 'TODO_ADD_SECOND_KEY'.
% 2. Inspect the embedded .bib database at the very top of this file to find the valid keys.
% 3. Correct the keys to match: huyen2024ai and kottwitz2021latex.
% 4. Run BibTeX (F8 in TeXstudio), followed by pdfLaTeX (F5) to resolve references.
% =========================================================================
Foundation model architectures and automated software agents require structured evaluation benchmarks \cite{huyen2099_wrong_key}. Concurrently, reproducible research documentation demands strict separation of content from typography through compiled markup systems \cite{TODO_ADD_SECOND_KEY}.

% =========================================================================
% SECTION 2: WORK BREAKDOWN STRUCTURE (WBS)
% =========================================================================
\section{Work Package Breakdown}

Project deliverables must be decomposed into explicit Work Packages (WPs). In Table~\ref{tab:wrong_label}, we summarize the primary delivery stages.

% =========================================================================
% [TODO - TASK 3]: BROKEN BOOKTABS TABLE & CROSS-REFERENCE
% Architectural Why: Academic tables avoid vertical rules and utilize booktabs
% horizontal rules for visual hierarchy. In tabularx, every row must contain
% exactly (N - 1) ampersands '&' for N columns, or the parser fails with
% an 'Extra alignment tab' or syntax error. Also, cross-references must point
% to an existing \label{} to resolve cleanly.
% Action:
% 1. Fix the missing column delimiter '&' in the WP2 row below.
% 2. Uncomment the \bottomrule command at the end of the tabular environment.
% 3. Fix the label/ref mismatch: Change '\ref{tab:wrong_label}' in the
%    text above to match the table's label '\label{tab:wp}'.
% =========================================================================
\begin{table}[htbp]
\centering\small
\renewcommand{\arraystretch}{1.3}
\begin{tabularx}{\textwidth}{l l X l l}
\toprule
\textbf{WP} & \textbf{Work Package} & \textbf{Main Tasks} & \textbf{Output} & \textbf{Weeks} \\
\midrule
WP1 & Research \& Ethics     & Literature review; formal ethics application & Approved Ethics & 1--4 \\
% ERROR BELOW: Notice the missing '&' between 'Requirements; system design' and 'Architecture Doc'
WP2 & Requirements & Requirements; system design Architecture Doc & 5--8 \\
WP3 & Core Implementation   & Artefact sprints; GitHub commits             & Working Artefact & 9--16 \\
% [TODO]: Uncomment the bottomrule line below to close the table properly
% \bottomrule
\end{tabularx}
\caption{Initial Work Package plan for Final Year Project}
\label{tab:wp}
\end{table}

% =========================================================================
% SECTION 3: 24-WEEK GANTT ROADMAP
% =========================================================================
\section{Project Timeline Schedule}

Visual scheduling in LaTeX is executed natively using the \texttt{pgfgantt} package, producing vector graphics that preserve typography and line weight without rasterization blur.

% =========================================================================
% [TODO - TASK 4]: INCOMPLETE PGFGANTT 24-WEEK SCHEDULE
% Architectural Why: pgfgantt renders tasks via coordinate units (weeks 1 to 24).
% \ganttgroup defines parent deliverables, \ganttbar defines task durations,
% and \ganttmilestone marks critical gating events (zero-duration checkpoints).
% Action:
% Complete WP2 (Weeks 5 to 8) and WP3 (Weeks 9 to 16) by uncommenting and
% completing the coordinate blocks below.
% =========================================================================
\begin{figure}[htbp]
\centering
\begin{ganttchart}[
    hgrid, vgrid,
    x unit=0.48cm,
    y unit chart=0.6cm,
    title label font=\tiny\bfseries,
    bar label font=\scriptsize,
    group label font=\scriptsize\bfseries,
    milestone label font=\scriptsize\itshape,
    bar/.append style={fill=NavyBlue!70},
    group/.append style={fill=Mahogany!80},
    milestone/.append style={fill=Goldenrod!90}
  ]{1}{24}
  \gantttitle{24-Week Final Year Project Timeline}{24} \\
  \gantttitlelist{1,...,24}{1} \\

  % --- Work Package 1 (Completed Reference) ---
  \ganttgroup{WP1: Research \& Ethics}{1}{4} \\
  \ganttbar{Literature Review}{1}{3} \\
  \ganttbar{Ethics Submission}{2}{4} \\
  \ganttmilestone{Ethics Approval Gate}{4} \\

  % --- Work Package 2: Requirements & Design (Weeks 5-8) ---
  % [TODO - TASK 4.1]: Uncomment and verify the coordinates for WP2
  % \ganttgroup{WP2: Requirements \& Design}{5}{8} \\
  % \ganttbar{System Specification}{5}{6} \\
  % \ganttbar{Architecture Modeling}{7}{8} \\
  % \ganttmilestone{Design Baseline}{8} \\

  % --- Work Package 3: Implementation (Weeks 9-16) ---
  % [TODO - TASK 4.2]: Complete WP3 implementation bar spanning Weeks 9 to 16
  % \ganttgroup{WP3: Core Implementation}{9}{16} \\
  % \ganttbar{Sprint 1--4 Development}{9}{16} \\
  % \ganttmilestone{Mid-Point Prototype}{12}

\end{ganttchart}
\caption{24-Week Project Roadmap with Gantt visualization}
\label{fig:gantt_sandbox}
\end{figure}

% =========================================================================
% REFERENCES (IEEE Numeric Style)
% =========================================================================
\bibliographystyle{ieeetr}
\bibliography{\jobname}

\end{document}
